\appendix

\section*{附录 A：核心概念速查表}

\begin{center}
\small
\begin{tabular}{>{\raggedright\arraybackslash}p{3.2cm} p{12.0cm}}
\toprule
\textbf{概念} & \textbf{含义} \\
\midrule
概念一 & 用一句话写清博主的定义，保留他本人的用词 \\
概念二 & …… \\
\bottomrule
\end{tabular}
\end{center}

\clearpage
\section*{附录 B：出处对照}

本讲义的思想内核提炼自博主的 N 期视频。下表为主要观点与原视频的对照，方便回看原作：

\begin{small}
\begin{longtable}{>{\raggedright\arraybackslash}p{5.2cm} >{\raggedright\arraybackslash}p{10.0cm}}
\toprule
\textbf{视频（标题）} & \textbf{本讲义中的相关章节 / 概念} \\
\midrule
\endfirsthead
\toprule
\textbf{视频（标题）} & \textbf{相关章节 / 概念（续）} \\
\midrule
\endhead
\bottomrule
\endfoot
示例标题 & 第 X 章某节、某概念 \\
\end{longtable}
\end{small}

\vspace{1em}

\term{观点属于博主，编排属于整理者；如有曲解，责任在整理一方。} 请把这份讲义当作一张地图，而不是一份结论——地图的作用是让你看清地形，路终究要自己走。
